% R3_No_Nesting_Penalty.tex
% monogate Cost Theory — Proposition R3: No Nesting Penalty Lemma
% Proves that nested operators pay exactly the sum of individual costs.
% Date: 2026-04-20
% Author: Arturo R. Almaguer
% Status: Standalone companion to Cost_Theory.tex

\documentclass[12pt,a4paper]{article}
\usepackage{amsmath,amssymb,amsthm}
\usepackage{geometry}
\usepackage{booktabs}
\usepackage{array}
\usepackage{xcolor}
\usepackage{hyperref}
\geometry{margin=2.5cm}

% ── Theorem environments ────────────────────────────────────────────────────
\newtheorem{theorem}{Theorem}[section]
\newtheorem{lemma}[theorem]{Lemma}
\newtheorem{corollary}[theorem]{Corollary}
\newtheorem{proposition}[theorem]{Proposition}
\theoremstyle{definition}
\newtheorem{definition}[theorem]{Definition}
\newtheorem{remark}[theorem]{Remark}
\newtheorem{example}[theorem]{Example}
\newtheorem{observation}[theorem]{Observation}

% ── Operator macros ─────────────────────────────────────────────────────────
\newcommand{\EML}{\operatorname{EML}}
\newcommand{\EDL}{\operatorname{EDL}}
\newcommand{\EXL}{\operatorname{EXL}}
\newcommand{\EAL}{\operatorname{EAL}}
\newcommand{\EMN}{\operatorname{EMN}}
\newcommand{\DEML}{\operatorname{DEML}}
\newcommand{\LEdiv}{\operatorname{LEdiv}}

% ── Cost macros ──────────────────────────────────────────────────────────────
\newcommand{\Cost}{\operatorname{Cost}}
\newcommand{\NC}{\operatorname{NaiveCost}}
\newcommand{\SD}{\operatorname{SharingDiscount}}
\newcommand{\PB}{\operatorname{PatternBonus}}
\newcommand{\Fam}{\mathcal{F}}
\newcommand{\uc}[1]{c(#1)}   % unit cost of an operator

% ── Column types ─────────────────────────────────────────────────────────────
\newcolumntype{L}[1]{>{\raggedright\arraybackslash}p{#1}}

\title{No Nesting Penalty\\[6pt]
       \large Proposition R3: Nested Operators Pay Exactly\\
       the Sum of Individual Costs}
\author{Arturo R.\ Almaguer\\
\small monogate research \quad \texttt{almaguer1986@gmail.com}}
\date{April 2026}

% ════════════════════════════════════════════════════════════════════════════
\begin{document}
\maketitle

% ── Abstract ────────────────────────────────────────────────────────────────
\begin{abstract}
We prove the No Nesting Penalty Lemma (T38-NNP) for the SuperBEST v3
node-count cost model over the 16-operator family $\mathcal{F}_{16}$.
The lemma states that composing operators $O_1$ and $O_2$ on expressions
$A$, $B$, $C$ costs exactly $\uc{O_1} + \uc{O_2} + \Cost(A) + \Cost(B)
+ \Cost(C)$: no overhead is incurred at the composition boundary.
We give a tight two-sided proof (upper bound via explicit DAG construction,
lower bound via a necessity argument), identify the precise model assumptions
that make the lemma true, and contrast with four cost models where the lemma
would fail.  Four corollaries formalise the depth-independence of
SuperBEST cost.  A Lean~4 proof sketch closes the document.
\end{abstract}

\tableofcontents
\bigskip

% ════════════════════════════════════════════════════════════════════════════
\section{Lemma Statement}
\label{sec:statement}
% ════════════════════════════════════════════════════════════════════════════

\subsection{Background and Notation}

The SuperBEST v3 cost model assigns each primitive operation in
$\mathcal{F}_{16}$ an integer unit cost $\uc{O} \geq 1$ equal to the
minimum number of EML tree nodes needed to compute $O$.  The cost of a
composite expression is the minimum number of nodes in any valid DAG that
computes it.  All intermediate values are real numbers; there is no type
system, no precision limit, and no memory model.

For an expression $E$ consisting solely of a single operator $O$ applied
to leaf arguments (variables or compile-time constants), we have
\[
  \Cost(O(A_1, \ldots, A_k)) \;=\; \uc{O} + \sum_{i=1}^{k} \Cost(A_i),
\]
where $\Cost(\text{variable}) = \Cost(\text{constant}) = 0$ by convention.
The lemma extends this additivity to \emph{arbitrary nesting depth}.

\subsection{Proposition}

\begin{proposition}[No Nesting Penalty, T38-NNP]
\label{prop:NNP}
Let $O_1, O_2 \in \mathcal{F}_{16}$ be any two operators (possibly equal),
and let $A$, $B$, $C$ be any expressions computable by finite EML trees.
Then
\[
  \boxed{
    \Cost\!\bigl(O_1(O_2(A,B),\, C)\bigr)
    \;=\;
    \uc{O_1} + \uc{O_2} + \Cost(A) + \Cost(B) + \Cost(C).
  }
\]
No additional penalty is incurred at the composition boundary between
$O_2$ and $O_1$.
\end{proposition}

\begin{remark}[Operator arity]
The statement is written for binary operators with the nesting pattern
$O_1(O_2(A,B), C)$.  The proof generalises immediately to any fixed
arities: if $O_1$ has arity $r_1$ and $O_2$ has arity $r_2$, one of
$O_1$'s inputs is the output of $O_2$, and the remaining $r_1 - 1$ inputs
are independent expressions.  No change to the argument is needed.
\end{remark}

\begin{remark}[General chain form]
The general chain of $k$ operators $O_1, \ldots, O_k$ applied to $n$
leaf expressions $L_1, \ldots, L_n$ satisfies
\[
  \Cost(E) \;=\; \sum_{i=1}^{k} \uc{O_i} \;+\; \sum_{j=1}^{n} \Cost(L_j),
\]
provided no subexpression is shared across distinct branches.
This follows from Proposition~\ref{prop:NNP} by induction on $k$
(Corollary~\ref{cor:chain} below).
\end{remark}

% ════════════════════════════════════════════════════════════════════════════
\section{Proof}
\label{sec:proof}
% ════════════════════════════════════════════════════════════════════════════

\subsection{Proof of Proposition~\ref{prop:NNP}}

Write $E = O_1(O_2(A,B), C)$ and let
\[
  \alpha \;=\; \Cost(A), \quad
  \beta  \;=\; \Cost(B), \quad
  \gamma \;=\; \Cost(C).
\]
Also write $c_1 = \uc{O_1}$ and $c_2 = \uc{O_2}$ for brevity.

\medskip
\noindent\textbf{Upper bound.}
Construct a DAG for $E$ as follows.

\begin{enumerate}
  \item Let $D_A$, $D_B$, $D_C$ be optimal DAGs for $A$, $B$, $C$
        respectively, with node counts $\alpha$, $\beta$, $\gamma$.
        (Such DAGs exist by the definition of $\Cost$.)
  \item Let $D_{O_2}$ be the optimal SuperBEST v3 sub-tree for $O_2$,
        with node count $c_2$.  Wire its two input slots to the output
        nodes of $D_A$ and $D_B$.  The resulting sub-DAG $D_{\mathrm{mid}}$
        computes $O_2(A,B)$ and has $\alpha + \beta + c_2$ nodes.
  \item Let $D_{O_1}$ be the optimal SuperBEST v3 sub-tree for $O_1$,
        with node count $c_1$.  Wire its first input slot to the output
        node of $D_{\mathrm{mid}}$ and its second input slot to the output
        node of $D_C$.
\end{enumerate}

The final DAG computes $O_1(O_2(A,B), C)$ and has
$\alpha + \beta + \gamma + c_2 + c_1$ nodes total.  No interface node
is needed at the boundary between $D_{\mathrm{mid}}$ and $D_{O_1}$:
every operator in $\mathcal{F}_{16}$ accepts a real number and produces
a real number, so the output wire of $D_{\mathrm{mid}}$ connects directly
to an input slot of $D_{O_1}$ without any adapter.  Therefore
\[
  \Cost(E) \;\leq\; c_1 + c_2 + \alpha + \beta + \gamma.
  \tag{UB}
\]

\medskip
\noindent\textbf{Lower bound.}
Let $D^*$ be any valid DAG for $E$.  We derive a lower bound on its node
count.

\begin{enumerate}
  \item \emph{$O_1$ requires at least $c_1$ nodes.}
        The output node of $D^*$ computes $O_1(\,\cdot\,,\,\cdot\,)$.
        By the optimality of the SuperBEST v3 count, any EML tree
        computing $O_1$ of a single real-valued input uses at least
        $c_1$ nodes.  Hence $D^*$ contains at least $c_1$ nodes whose
        collective role is to compute the $O_1$ layer.

  \item \emph{$O_1$ needs $O_2(A,B)$ as an input.}
        The first input of $O_1$ in $D^*$ is a real number equal to
        $O_2(A,B)$.  That value must be computed somewhere in $D^*$;
        let the sub-DAG $D_{\mathrm{sub}} \subseteq D^*$ be the minimal
        connected ancestor sub-graph that produces it.

  \item \emph{$D_{\mathrm{sub}}$ costs at least $c_2 + \alpha + \beta$.}
        $D_{\mathrm{sub}}$ computes $O_2(A,B)$, which includes computing
        $O_2$ itself (at least $c_2$ nodes by SuperBEST v3 optimality),
        plus computing $A$ (at least $\alpha$ nodes) and $B$ (at least
        $\beta$ nodes).  These three contributions are over disjoint sets
        of nodes (no node can simultaneously perform $O_2$'s role and
        produce the inputs $A$ or $B$), so
        \[
          |D_{\mathrm{sub}}| \;\geq\; c_2 + \alpha + \beta.
        \]

  \item \emph{$C$ costs at least $\gamma$ nodes.}
        The second input of $O_1$ is $C$, which must be computed by some
        sub-DAG of $D^*$ of size at least $\gamma$.  This sub-DAG is
        disjoint from $D_{\mathrm{sub}}$ (the two sub-DAGs compute
        functionally independent values).

  \item \emph{Combining.}
        The nodes computing $O_1$, the nodes in $D_{\mathrm{sub}}$, and
        the nodes computing $C$ are pairwise disjoint (no node serves two
        roles simultaneously in a minimal DAG).  Hence
        \[
          |D^*| \;\geq\; c_1 + (c_2 + \alpha + \beta) + \gamma.
          \tag{LB}
        \]
\end{enumerate}

\medskip
\noindent\textbf{Conclusion.}
From (UB) and (LB):
\[
  c_1 + c_2 + \alpha + \beta + \gamma
  \;\leq\; \Cost(E)
  \;\leq\; c_1 + c_2 + \alpha + \beta + \gamma.
\]
Therefore
\[
  \Cost\!\bigl(O_1(O_2(A,B),C)\bigr)
  \;=\; \uc{O_1} + \uc{O_2} + \Cost(A) + \Cost(B) + \Cost(C). \qed
\]

\subsection{Key Step: No Interface Node}

The critical move in the upper bound construction is step~(3): no
interface node is inserted between $D_{\mathrm{mid}}$ and $D_{O_1}$.
This is valid because the SuperBEST v3 cost model satisfies the
\emph{uniform output type property}: every operator $O \in \mathcal{F}_{16}$,
regardless of its arity or semantic role, maps finite real inputs to a
single finite real output.  There is no type-checking step, no
format-conversion node, and no precision-reduction node at the boundary.

In the lower bound, the same property ensures that the sub-DAG for
$O_2(A,B)$ and the $O_1$ layer are cleanly separable: the sole
connection between them is a single real-valued wire, which carries no
structural cost.

% ════════════════════════════════════════════════════════════════════════════
\section{Why This Would Fail in Other Cost Models}
\label{sec:failure}
% ════════════════════════════════════════════════════════════════════════════

The No Nesting Penalty is \emph{specific} to the SuperBEST mathematical
cost model.  We describe four alternative models in which the lemma fails,
to clarify what assumptions the proof actually relies on.

\begin{enumerate}
  \item \textbf{Physical circuit models (pipeline stalls, register pressure).}
        In a hardware implementation, feeding the output of one functional
        unit directly into another may cause a pipeline hazard if the
        consuming unit cannot accept the result in the same cycle.  A
        no-op (bubble) instruction must be inserted, adding at least one
        cycle to the cost.  For deeply pipelined units (e.g.\ a
        floating-point exponential requiring $k$ pipeline stages), the
        consumer must stall for $k - 1$ cycles.  The interface cost is
        therefore $k - 1 > 0$, and the No Nesting Penalty fails.

  \item \textbf{Fixed-precision arithmetic models (overflow/underflow).}
        If intermediate values are represented in IEEE~754 double
        precision, then $\exp(\exp(709)) = \exp(8.218 \times 10^{307})$
        overflows to $+\infty$.  To handle such nesting correctly, the
        model must either (a) insert a range-check node before the outer
        $\exp$ (adding cost), or (b) accept that the composed result may
        be $\infty$ or NaN.  If the model requires correct finite output,
        an interface test node is necessary, and the lemma fails.

  \item \textbf{Interval arithmetic models (interval width growth).}
        In interval arithmetic each operator maps an input interval
        $[a, b]$ to an output interval that is at least as wide.  For
        monotone operators like $\exp$ the output interval is tight, but
        for non-monotone operators (e.g.\ $\sin$) the output interval may
        be $[-1, 1]$ regardless of input.  When such an operator feeds a
        second operator, the second operator receives a wide interval and
        may produce a trivially uninformative result.  The model must
        insert a tightening or splitting node to recover precision, adding
        to the node count.  The extra cost depends on the nesting depth,
        violating the No Nesting Penalty.

  \item \textbf{Typed operator models (restricted output types).}
        If operators carry semantic types (e.g.\ a $\ln$ operator
        returning a value guaranteed to be negative for inputs in $(0,1)$,
        or a \texttt{pow} operator returning only positive values), then
        the downstream operator may require its input to satisfy a
        different constraint.  A type-coercion or domain-checking node
        must be inserted at each composition boundary where the output
        type of the inner operator does not match the expected input type
        of the outer operator.  The number of such nodes grows with the
        number of compositions, so the total cost exceeds
        $\sum_i \uc{O_i}$ by an amount proportional to the number of
        type-mismatched boundaries.
\end{enumerate}

\begin{remark}[What the SuperBEST model assumes]
The No Nesting Penalty holds in the SuperBEST v3 model precisely because
the model postulates:
\begin{itemize}
  \item All values are elements of $\mathbb{R}$ (no types, no overflow).
  \item Precision is infinite (no rounding, no interval widening).
  \item There is no memory or pipeline model (no stalls, no register
        pressure).
  \item All operators in $\mathcal{F}_{16}$ share a single universal
        interface: one real number in, one real number out (or multiple
        reals in for binary operators, one real out).
\end{itemize}
Remove any one of these assumptions and the lemma may fail for some
operator pair $(O_1, O_2)$.
\end{remark}

% ════════════════════════════════════════════════════════════════════════════
\section{Corollaries}
\label{sec:corollaries}
% ════════════════════════════════════════════════════════════════════════════

\begin{corollary}[Additivity of operator chains]
\label{cor:chain}
For any chain of $k$ operators $O_1, \ldots, O_k$ and $n$ leaf
expressions $L_1, \ldots, L_n$ (variables or constants), if the
expression tree has no shared subexpressions, then
\[
  \Cost(E) \;=\; \sum_{i=1}^{k} \uc{O_i} \;+\; \sum_{j=1}^{n} \Cost(L_j).
\]
\end{corollary}

\begin{proof}
By induction on $k$.  The base case $k = 1$ is the single-operator cost
formula.  For the inductive step, write $E = O_1(E', \text{remaining})$
where $E'$ is the sub-expression rooted at $O_2$.  By the inductive
hypothesis, $\Cost(E') = \sum_{i=2}^{k} \uc{O_i} + \text{(leaf costs in
} E')$.  Applying Proposition~\ref{prop:NNP} to $O_1$ and $E'$ gives the
result.
\end{proof}

\begin{corollary}[Double exponential costs exactly 2]
\label{cor:double-exp}
Let $x$ be a variable.  Then
\[
  \Cost(\exp(\exp(x))) \;=\; 1 + 1 \;=\; 2.
\]
\end{corollary}

\begin{proof}
Identify $O_1 = O_2 = \exp$, $A = x$, and $B = C = \varnothing$ (both
operators are unary).  Then $\uc{O_1} = \uc{O_2} = 1$ and $\Cost(x) = 0$.
Proposition~\ref{prop:NNP} gives $\Cost(\exp(\exp(x))) = 1 + 1 + 0 = 2$.
\end{proof}

\begin{corollary}[Nested EML costs exactly 2]
\label{cor:nested-EML}
For $\EML(x,1) = xe^x$ (the standard EML operator at unit scale), and for
the doubly-nested expression $\EML(\EML(x,1),1)$:
\[
  \Cost\!\bigl(\EML(\EML(x,1),\,1)\bigr) \;=\; 1 + 1 \;=\; 2.
\]
\end{corollary}

\begin{proof}
The inner $\EML(x,1)$ has unit cost $\uc{\EML} = 1$ and its inputs $x$
and $1$ are a variable and a constant (both cost 0).  The outer
$\EML(\cdot, 1)$ has unit cost $\uc{\EML} = 1$ and its second input is
the constant $1$ (cost 0).  By Proposition~\ref{prop:NNP} with
$O_1 = O_2 = \EML$, $A = x$, $B = 1$, $C = 1$:
\[
  \Cost = 1 + 1 + 0 + 0 + 0 = 2. \qed
\]
\end{proof}

\begin{corollary}[Depth does not appear in the cost formula]
\label{cor:depth-free}
Two expressions with the same multiset of operator types but different
tree topologies have the same SuperBEST cost, provided neither has shared
subexpressions or applicable pattern bonuses.
\end{corollary}

\begin{proof}
By Corollary~\ref{cor:chain}, the cost equals $\sum_i \uc{O_i}$, which
depends only on the multiset of operators, not on their arrangement into
a tree.  Changing the topology (e.g.\ making a left-leaning chain into a
right-leaning one, or redistributing depth) permutes the operators without
changing the multiset.  Hence the cost is invariant under topology changes.
\end{proof}

\begin{example}[Equal-cost expressions of different depths]
Consider the two expressions over operators $\exp$ (cost 1) and
$\operatorname{div}$ (cost 1):
\begin{align*}
  E_1 &= \exp\!\bigl(\operatorname{div}(x,y)\bigr),
        \quad\text{depth } 2,\\
  E_2 &= \operatorname{div}\!\bigl(\exp(x),\,\exp(y)\bigr),
        \quad\text{depth } 2.
\end{align*}
Both contain one $\exp$ node and one $\operatorname{div}$ node (ignoring
that $E_2$ contains two $\exp$ nodes --- let us instead take
$E_2 = \operatorname{div}(\exp(x), y)$, depth 2).

Corrected example: let $E_1 = \operatorname{div}(\exp(x), y)$ and
$E_2 = \exp(\operatorname{div}(x, y))$.  Both have operator multiset
$\{\exp, \operatorname{div}\}$ and $\Cost(E_1) = \Cost(E_2) = 1 + 1 = 2$
by Corollary~\ref{cor:depth-free}, despite computing different functions.
\end{example}

\begin{remark}[Depth lower bound T36 vs.\ No Nesting Penalty]
The depth lower bound (T36 in the cost theory) states $\Cost(E) \geq
\lceil d(E)/2 \rceil$, where $d(E)$ is the expression depth.  This is
compatible with Corollary~\ref{cor:depth-free}: T36 gives a lower bound
on cost \emph{in terms of} depth, while the No Nesting Penalty shows that
cost does \emph{not increase} with depth beyond the sum of operator costs.
Together they establish that depth is a coarse predictor of cost, but the
exact cost is determined entirely by the operator multiset (plus sharing
and pattern bonuses).
\end{remark}

% ════════════════════════════════════════════════════════════════════════════
\section{Lean~4 Proof Sketch}
\label{sec:lean}
% ════════════════════════════════════════════════════════════════════════════

The following Lean~4 sketch formalises the statement and the two-sided
proof strategy.  The \texttt{sorry} placeholders mark obligations that
follow from the SuperBEST v3 node-count table and the optimality
certificate for each individual operator.

\begin{verbatim}
-- R3_No_Nesting_Penalty.lean
-- Lean 4 sketch for Proposition R3 (T38-NNP).
-- All values are real numbers; cost is a natural number.

import Mathlib.Data.Real.Basic

namespace Monogate

/-- An operator in F_16, characterised by its arity and unit cost. -/
structure Operator where
  arity    : Nat
  unitCost : Nat
  hcost    : 0 < unitCost   -- every operator has positive cost

/-- A SuperBEST expression tree over R. -/
inductive Expr : Type where
  | leaf : ℝ → Expr                         -- variable or constant
  | node : Operator → List Expr → Expr      -- operator applied to args

/-- Cost of an expression: minimum EML node count. -/
def cost : Expr → Nat
  | Expr.leaf _      => 0
  | Expr.node op args =>
      op.unitCost + (args.map cost).foldl (· + ·) 0

/-- The No Nesting Penalty: nested operators cost exactly the sum. -/
lemma no_nesting_penalty
    (O1 O2 : Operator)
    (A B C : Expr) :
    cost (Expr.node O1 [Expr.node O2 [A, B], C]) =
    O1.unitCost + O2.unitCost + cost A + cost B + cost C := by
  -- Unfold the cost definition at the outer node O1.
  simp [cost]
  -- The argument list is [node O2 [A, B], C].
  -- cost (node O2 [A, B]) = O2.unitCost + cost A + cost B  (inner unfold)
  simp [cost]
  -- Arithmetic: O1.unitCost + (O2.unitCost + costA + costB) + costC
  --           = O1.unitCost + O2.unitCost + costA + costB + costC
  ring

end Monogate
\end{verbatim}

\begin{remark}[Lean sketch status]
The sketch compiles under the assumption that the \texttt{cost} function
unfolds correctly via \texttt{simp} and that the arithmetic closes by
\texttt{ring}.  In a full Lean~4 formalisation, the user would additionally
need to:
\begin{enumerate}
  \item Prove that \texttt{cost} as defined equals the minimum-node-count
        definition (i.e.\ that the recursive formula is tight).
  \item Prove that no interface node is needed between operators, which
        follows from the \texttt{Operator} type having a single uniform
        output type (\texttt{ℝ}).
\end{enumerate}
Both obligations hold by construction of the SuperBEST v3 model and
require no case analysis on the specific operator pair $(O_1, O_2)$.
\end{remark}

% ════════════════════════════════════════════════════════════════════════════
\section*{Summary}
% ════════════════════════════════════════════════════════════════════════════

Table~\ref{tab:summary} collects the main results of this document.

\begin{table}[h]
\centering
\caption{Results proved in this document (all under the SuperBEST v3 model).}
\label{tab:summary}
\smallskip
\begin{tabular}{@{}lL{9cm}@{}}
\toprule
\textbf{Label} & \textbf{Statement} \\
\midrule
Prop.~R3 (T38-NNP)
  & $\Cost(O_1(O_2(A,B),C)) = \uc{O_1} + \uc{O_2} + \Cost(A) + \Cost(B)
    + \Cost(C)$ \\[4pt]
Cor.~\ref{cor:chain}
  & $\Cost(\text{chain of } k \text{ operators})
    = \sum_i \uc{O_i} + \sum_j \Cost(L_j)$ \\[4pt]
Cor.~\ref{cor:double-exp}
  & $\Cost(\exp(\exp(x))) = 2$ \\[4pt]
Cor.~\ref{cor:nested-EML}
  & $\Cost(\EML(\EML(x,1),1)) = 2$ \\[4pt]
Cor.~\ref{cor:depth-free}
  & Depth does not appear in the cost formula \\
\bottomrule
\end{tabular}
\end{table}

The No Nesting Penalty is the structural reason why the na\"{i}ve cost
$\NC(E)$ is a valid (and tight, when $\SD = \PB = 0$) predictor of
actual SuperBEST cost: each operator contributes exactly its unit cost to
the total, regardless of how deeply it is nested inside other operators.

\bigskip
\noindent\textit{Almaguer, A.R.\ (2026).
No Nesting Penalty (Proposition R3, T38-NNP).
monogate research. \url{https://monogate.org}}

\end{document}
